\documentclass[10pt]{article}
\usepackage[margin=1in]{geometry}
\usepackage{amsmath,amssymb,booktabs,array,microtype}
\usepackage[hidelinks]{hyperref}
\usepackage{enumitem}
\usepackage{graphicx}
\usepackage{siunitx}

\title{Recoverable Feature Routing under Distribution Shift:\\
A Controlled Study of Gradient Accessibility and Selective Fast Adaptation}
\author{Anonymous draft}
\date{October 2026}

\begin{document}
\maketitle

\begin{abstract}
Neural classifiers commonly prioritize features that are most predictive on the training distribution. This is efficient when feature reliability is stable, but brittle when a dominant feature fails while lower-priority evidence remains informative. We study this failure mode in a controlled synthetic image setting containing three explicitly isolated evidence families: a dominant but fragile eye-region cue, a weaker lower-face geometric cue, and a weaker chromatic cue. A held-out glasses perturbation invalidates the dominant cue without removing the two backup cues. Across five seeds, fixed global fusion falls to 0.5000 accuracy under this shift, while a few-shot adaptive router recovers to $0.6302 \pm 0.1064$ at 8 shots. Recovery is seed-dependent because some routers become effectively locked to the failed feature. Anti-saturation transformations reduce apparent output collapse but do not consistently restore adaptation, showing that nonzero routing mass is not equivalent to gradient accessibility. An independent fast routing path raises held-out accuracy to $0.6986 \pm 0.0259$ at 32 shots, but applying the fast policy globally reduces clean accuracy to 0.7158. We therefore introduce support-keyed selective fast-path activation. A linear memory switch trained only to distinguish the new regime from a small clean reference memory matches the fast-only shifted accuracy while retaining 1.0000 clean accuracy: $0.6986 \pm 0.0259$ on glasses and $1.0000 \pm 0.0000$ on clean data at 32 shots, with zero clean false activations in this controlled setting. The result separates three requirements for adaptive evidence use: retaining backup features, preserving gradient-accessible routes to them, and selectively activating those routes under shift. We call this combination \emph{recoverable feature routing}.
\end{abstract}

\section{Introduction}

Feature selection is usually discussed as a question of global utility: features that contribute strongly across a dataset are preserved or amplified, while weak features are downweighted, pruned, or ignored. This view is reasonable when feature reliability is approximately stationary. It is less satisfactory when a feature that is weak on average becomes useful for a specific input regime, or when the feature that is normally strongest becomes unreliable.

The motivating example is deepfake detection. A detector may learn to rely strongly on a particular family of texture, frequency, chromatic, geometric, or localized facial cues. An intervention such as eyewear, occlusion, compression, or a generator-specific artifact can alter the reliability of one family without destroying all remaining evidence. In such a case, the desired behavior is not necessarily permanent retraining. The model should be able to temporarily reconsider which evidence it trusts.

This paper isolates that mechanism in a deliberately controlled synthetic setting before attempting a full deepfake benchmark. The experiments are built around three questions:
\begin{enumerate}[leftmargin=*]
    \item Can globally weak evidence remain useful after the dominant feature fails?
    \item If useful backup evidence still exists, can an adaptive router recover access to it?
    \item If a fully independent fast path restores adaptation but damages old-distribution performance, can that path be invoked selectively?
\end{enumerate}

The central distinction that emerges is between \emph{feature retention} and \emph{feature recoverability}. A representation may still exist in the model and may even retain nonzero routing probability, yet remain functionally inaccessible to rapid adaptation because the routing map has entered a low-gradient state.

\section{Related Work}

\subsection{Conditional computation and mixture-of-experts routing}

Mixture-of-experts models use input-dependent gates to select or combine specialized subnetworks. Sparsely gated MoE layers demonstrated that conditional computation can increase model capacity without activating all parameters for every input~\cite{shazeer2017}. Switch Transformers later simplified sparse routing and highlighted training instabilities associated with learned expert selection~\cite{fedus2022}. Our setting is much smaller and uses routing as an instrument for studying recoverability rather than capacity scaling.

\subsection{Few-shot and fast adaptation}

Model-Agnostic Meta-Learning explicitly optimizes a model so that a small number of gradient steps on a new task can produce useful adaptation~\cite{finn2017}. Fast-weight approaches distinguish rapidly changing variables from slower learned parameters~\cite{ba2016}. Our experiments adopt this slow/fast distinction but focus on a different failure mode: a fast update can itself be ineffective if its gradient must pass through a saturated slow routing map.

\subsection{Test-time adaptation and stability-plasticity}

Tent adapts model parameters at inference time under distribution shift~\cite{wang2021}. Continual-learning work studies the complementary problem of learning new behavior while retaining previous capabilities. Elastic Weight Consolidation preserves old-task behavior by constraining changes to important parameters~\cite{kirkpatrick2017}. Our setting exposes a local version of the same stability-plasticity tension.

\section{Controlled Experimental Setting}

\subsection{Evidence families}

Each synthetic image contains three separately accessible evidence families. \textbf{A} is a nearly deterministic eye-region cue and is dominant on clean data. \textbf{B} is a lower-face geometric cue correlated with the label with probability 0.72. \textbf{C} is a chromatic cue correlated with the label with probability 0.66.

Three small convolutional experts receive isolated feature views: the eye band, grayscale lower-face geometry, and opponent-color channels.

\subsection{Held-out intervention}

A synthetic glasses overlay is held out from slow-router training. It invalidates A while leaving B and C largely intact. The slow router is meta-trained on clean inputs, eye-region masking, lower-face masking, desaturation, blur, and JPEG compression.

\subsection{Evaluation}

The main runs use seeds 11, 22, 33, 44, and 55. Few-shot support sizes are 2, 4, 8, 16, and 32. Support samples are excluded from query evaluation.

\section{Experiments}

\subsection{Experiment 1: overly easy toy world}

The first construction used RGB, edge, high-pass, and chromatic experts. All fusion strategies achieved 1.000 accuracy on both clean and glasses-corrupted data. The perturbation increased cross-entropy but left the edge cue sufficiently predictive that no meaningful adaptive behavior was required.

\subsection{Experiment 2: fragile dominance and few-shot recovery}

The revised construction forced a fragile dominant cue. Fixed global fusion fell to 0.5000 mean glasses accuracy across five seeds. The adaptive router reached $0.6302 \pm 0.1064$ at 8 shots, while a support-optimized fusion reference reached $0.6944 \pm 0.0302$ at 32 shots.

\begin{table}[h]
\centering
\caption{Experiment 2 shifted-regime accuracy.}
\begin{tabular}{lrr}
\toprule
Method & $K$ & Accuracy \\
\midrule
Global fixed fusion & 0 & $0.5000 \pm 0.0000$ \\
Adaptive zero-shot & 0 & $0.5008 \pm 0.0016$ \\
Adaptive few-shot & 2 & $0.5872 \pm 0.1052$ \\
Adaptive few-shot & 8 & $\mathbf{0.6302 \pm 0.1064}$ \\
Adaptive few-shot & 32 & $0.6300 \pm 0.1061$ \\
Support-optimized fusion & 32 & $0.6944 \pm 0.0302$ \\
\bottomrule
\end{tabular}
\end{table}

At 8 shots, mean routing changed from approximately
\[
[A,B,C]=[0.904,0.088,0.008]
\]
to
\[
[A,B,C]=[0.451,0.539,0.010].
\]
However, the high variance revealed seed-dependent routing lock-in.

\subsection{Experiment 3: output anti-collapse is not gradient recoverability}

Temperature scaling, epsilon floors, logit clipping, and bounded residual routing were compared with standard softmax. The epsilon floor eliminated measured output saturation but reached only $0.5444 \pm 0.0878$ at 32 shots. Logit clipping peaked at $0.5877 \pm 0.1074$, similar to standard softmax.

For
\[
w_i=(1-E\epsilon)\operatorname{softmax}(z)_i+\epsilon,
\]
the output can be nonzero while the derivative with respect to $z_i$ remains controlled by the saturated softmax Jacobian. Thus nonzero routing mass is not equivalent to gradient access.

\subsection{Experiment 4: independent fast pathways}

Experiment 4 trained one shared slow router per seed and compared adaptation mechanisms on identical support/query splits.

The independent mixture uses
\[
p_s=\operatorname{softmax}(z_s), \qquad
p_f=\operatorname{softmax}(a),
\]
and
\[
w=(1-\alpha)p_s+\alpha p_f.
\]

\begin{table}[h]
\centering
\caption{Best operating point per Experiment 4 method.}
\begin{tabular}{lrrr}
\toprule
Method & $K$ & Glasses & Clean after \\
\midrule
Slow SGD & 8 & $0.6302 \pm 0.1064$ & 0.9226 \\
Additive fast & 8 & $0.6302 \pm 0.1064$ & 0.8668 \\
Fixed mixture & 32 & $0.5657 \pm 0.0805$ & 1.0000 \\
Learned mixture & 4 & $0.6070 \pm 0.0954$ & $\mathbf{0.9974}$ \\
Fast only & 32 & $\mathbf{0.6986 \pm 0.0259}$ & 0.7158 \\
\bottomrule
\end{tabular}
\end{table}

The independent fast route improved both mean shifted accuracy and seed stability. Its routing conductance at initialization was 0.4714, whereas the slow route was much smaller and approached zero on the most strongly locked seeds. The price of applying the fast route globally was a substantial clean-distribution degradation.

\subsection{Experiment 5: support-keyed selective activation}

Experiment 5 closes the stability-plasticity loop exposed by Experiment 4. The expert bank, slow router, support set, query set, and adapted fast policy are held fixed across activation methods. The only experimental variable is when the fast path is invoked.

The final route is
\[
w(x)=(1-\alpha(x))p_s(x)+\alpha(x)p_f.
\]

Five seeds are evaluated with $K\in\{2,4,8,16,32\}$. A balanced clean reference memory of 128 samples is drawn from the validation distribution and is never taken from the clean test set.

The prototype switch uses a frozen context embedding $\phi(x)$:
\[
s(x)=
\cos(\phi(x),\mu_{\mathrm{support}})
-
\cos(\phi(x),\mu_{\mathrm{clean}}).
\]

The memory switch instead trains a small linear environment classifier on frozen context features, using glasses support examples as the new-regime class and the clean reference memory as the old-regime class. It never receives fake/real targets.

The principal joint metric is
\[
H=
\frac{2A_{\mathrm{glasses}}A_{\mathrm{clean}}}
{A_{\mathrm{glasses}}+A_{\mathrm{clean}}}.
\]

\begin{table}[h]
\centering
\caption{Experiment 5 best operating point by joint stability-plasticity score.}
\begin{tabular}{lrrrr}
\toprule
Method & $K$ & Glasses & Clean & $H$ \\
\midrule
Slow SGD & 8 & $0.6302 \pm 0.1064$ & $0.9226 \pm 0.1081$ & 0.7367 \\
Fast only & 4 & $0.6669 \pm 0.0701$ & $0.8018 \pm 0.1270$ & 0.7179 \\
Global mixture & 8 & $0.5925 \pm 0.0836$ & $0.9674 \pm 0.0652$ & 0.7294 \\
Prototype switch & 32 & $0.5636 \pm 0.0790$ & $1.0000 \pm 0.0000$ & 0.7177 \\
\textbf{Memory switch} & \textbf{32} & $\mathbf{0.6986 \pm 0.0259}$ & $\mathbf{1.0000 \pm 0.0000}$ & \textbf{0.8223} \\
Environment oracle & 32 & $0.6986 \pm 0.0259$ & $1.0000 \pm 0.0000$ & 0.8223 \\
\bottomrule
\end{tabular}
\end{table}

The memory switch exactly matches the environment-oracle accuracies at every tested support size in this synthetic setting. At $K=32$,
\[
\bar{\alpha}_{\mathrm{glasses}}=1.000,
\qquad
\bar{\alpha}_{\mathrm{clean}}\approx0.0488,
\]
with a glasses activation rate of 1.000 and a clean false-activation rate of 0.000.

The prototype switch is also informative. Its thresholded regime decision is nearly perfect at $K=32$, but its mean mixture coefficients remain close to 0.5:
\[
\bar{\alpha}_{\mathrm{glasses}}\approx0.5105,
\qquad
\bar{\alpha}_{\mathrm{clean}}\approx0.4939.
\]
Thus the representation separates the regimes, but the continuous prototype-to-mixture calibration is too weak. This is a calibration failure rather than a failure of regime discrimination.

In the controlled toy setting, context-dependent fast-path activation almost completely removes the stability-plasticity trade-off observed in Experiment 4. This is a mechanism demonstration rather than evidence that the same separation will be easy on natural deepfakes.

\section{Interpretation}

The experiments suggest a hierarchy of requirements:
\begin{enumerate}[leftmargin=*]
    \item \textbf{Feature presence is insufficient.} B and C can remain predictive while the shifted model fails.
    \item \textbf{Nonzero routing weight is insufficient.} Probability floors can preserve mass without restoring useful gradients.
    \item \textbf{Fresh parameters are insufficient if they inherit the saturated map.} Additive fast logits can remain trapped.
    \item \textbf{Independent gradient access improves recoverability.} A separate fast distribution restores much of the lost adaptation.
    \item \textbf{Independent access alone is insufficient for stability.} Applying the fast route indiscriminately damages familiar inputs.
    \item \textbf{Selective activation can close the gap in the controlled setting.} The memory switch matches the environment-oracle accuracies while retaining perfect clean accuracy across the five tested seeds.
\end{enumerate}

We call the combined goal \emph{recoverable feature routing}: preserving evidence, preserving an adaptation path to that evidence, and selectively invoking that path when the slow policy becomes unreliable.

\section{Limitations}

The current study is synthetic; it does not establish deepfake-detection performance. The glasses intervention is deterministic, the expert views are explicitly factorized, and adaptation uses labeled support examples. The experiments consider one dominant-feature failure at a time. Experiment 5 additionally retains a small clean reference memory, which should be compared with source-free alternatives.

\section{Path to a Real Deepfake Study}

A real follow-up should use identity-disjoint splits, generator-disjoint evaluation, natural glasses and sunglasses, generated eyewear interventions, compression and blur controls, and multiple measurable evidence families such as spatial texture, frequency residuals, chromatic statistics, geometry, and physically motivated cues.

The target claim should remain mechanistic rather than benchmark-driven: determine whether a dominant cue becomes unreliable under a real condition, whether lower-priority cues remain informative, whether standard adaptation loses access to those cues, and whether an independent selective fast path restores them without degrading ordinary inputs.

\section{Conclusion}

A model can retain useful evidence while losing practical access to it through an overconfident routing policy. Output-level anti-collapse does not guarantee gradient accessibility. Independent fast routing substantially improves recovery under the controlled shift, but creates a stability-plasticity trade-off when applied globally.

Experiment 5 shows that this trade-off is not inevitable in the controlled setting. A small memory-based environment switch activates the fast path on glasses inputs while preserving the slow route on clean inputs. At 32 shots it reaches $0.6986 \pm 0.0259$ glasses accuracy and $1.0000 \pm 0.0000$ clean accuracy, matching the environment-oracle accuracies in this synthetic experiment.

The resulting hypothesis is more specific than dynamic feature weighting: adaptive evidence use requires \emph{retention, recoverability, and selective activation}. The next step is to test whether the same failure-and-recovery structure appears in realistic deepfake detection, where feature families are entangled and regime boundaries are substantially less obvious.

\begin{thebibliography}{9}

\bibitem{shazeer2017}
N. Shazeer, A. Mirhoseini, K. Maziarz, A. Davis, Q. V. Le, G. E. Hinton, and J. Dean,
``Outrageously Large Neural Networks: The Sparsely-Gated Mixture-of-Experts Layer,''
\emph{ICLR}, 2017.

\bibitem{fedus2022}
W. Fedus, B. Zoph, and N. Shazeer,
``Switch Transformers: Scaling to Trillion Parameter Models with Simple and Efficient Sparsity,''
\emph{Journal of Machine Learning Research}, vol. 23, no. 120, pp. 1--39, 2022.

\bibitem{finn2017}
C. Finn, P. Abbeel, and S. Levine,
``Model-Agnostic Meta-Learning for Fast Adaptation of Deep Networks,''
\emph{Proceedings of the 34th International Conference on Machine Learning},
PMLR 70, pp. 1126--1135, 2017.

\bibitem{ba2016}
J. Ba, G. Hinton, V. Mnih, J. Z. Leibo, and C. Ionescu,
``Using Fast Weights to Attend to the Recent Past,''
\emph{Advances in Neural Information Processing Systems}, 2016.

\bibitem{wang2021}
D. Wang, E. Shelhamer, S. Liu, B. Olshausen, and T. Darrell,
``Tent: Fully Test-Time Adaptation by Entropy Minimization,''
\emph{International Conference on Learning Representations}, 2021.

\bibitem{kirkpatrick2017}
J. Kirkpatrick et al.,
``Overcoming Catastrophic Forgetting in Neural Networks,''
\emph{Proceedings of the National Academy of Sciences},
vol. 114, no. 13, pp. 3521--3526, 2017.

\end{thebibliography}

\end{document}
